
\subsection*{Introduction}

For the last decades, Artificial Intelligence has become a powerful tool 
for enterprises but also more personally, especially as home assistant. 
The famous ones are Large Language Models such as ChatGPT, Gemini, and others 
because we can naturally chat with them.
The other famous type of models are classifiers and are still very used 
today. A classifier is a model (Random Forest, Neural Network, etc...) which 
sort inputs into \emph{classes}. For example, we can train a classifier to 
recognize cat images from dog ones or tumor scan from benign outgrowth scan.
The natural question is then "How to trust the decisions of a classifier?"
For non critic uses, we can simply test the model for several inputs 
to see if the model is able to recognize cat from dogs. But in sensitive 
domains, we may want to verify that the model effectively recognizes cat 
instead of wool ball or cat tree in background. In this case, the model 
could effectively make an error and class a tumor scan in the benign scan class and then 
create a medical error. 

When we use non-complex models such as little Random Forest, we can 
verify each decision of the model during the inference process but for 
more complex ones like Neural Networks with billions of neurons, the human
eye cannot do the same. 
That's why, researchers try to develop \emph{explanation methods} which aims 
to provide a human comprehensible output about the black-box model decision process 
for the user. The goal is to explain why the model took a certain decision 
and to figure out the origin of the potential errors. 

This document detailed the today used models and systems for classifier 
explanation. We start by presenting the SHAP-Score one which assign an importance 
value to each model inputs. This method is based on game theory. 
Then, we develop perturbations methods with LIME and Anchor systems. 
We present the theoretical aspects of these methods and their ameliorations. 
Finally, a more theoretical section develops Abdductive Explanation methods. 
This section tries to formalize a classifier's reasoning to provide a formal 
guaranty of the produced explanation. 

% \begin{table}[ht]
%     \centering
%     \caption{Comparative analysis of XAI methods (Perturbation vs. Abductive approaches).}
%     \label{tab:comparaison_xai}
%     \vspace{0.3cm}
    
%     \newcolumntype{L}{>{\RaggedRight\arraybackslash}X}
    
%     \small
%     \begin{tabularx}{\textwidth}{@{} l c c c c L @{}}
%         \toprule
%         \textbf{Method} & \textbf{Complexity} & \textbf{Approx.} & \textbf{Reach} & \textbf{Reliability} & \textbf{Key Characteristics} \\
%         \midrule
        
%         \textbf{LIME} & Polynomial & \checkmark & Local & Low & 
%         Uses a local linear surrogate. Fast but lacks formal guarantees and stability. \\
%         \addlinespace[2pt]
        
%         \textbf{SHAP} & Exponential & \checkmark & Local & Medium & 
%         Game theory based (Shapley values). Computationally heavy for high-dimensional data. \\
%         \addlinespace[2pt]
        
%         \textbf{Anchors} & Exponential & \checkmark & Local & Medium & 
%         IF-THEN rules as sufficient conditions. Combinatorial search space for large feature sets. \\
%         \midrule
        
%         \textbf{sbAXp} & Polynomial & \xmark & Local & \checkmark & 
%         Sample-based Abductive Explanations. Provides logical guarantees on a subset of samples. \\
%         \addlinespace[2pt]
        
%         \textbf{Argument} & Polynomial & \xmark & Local & \checkmark & 
%         Uses argumentation theory to derive explanations from conflicting data points. \\
%         \addlinespace[2pt]
        
%         \textbf{CPI-Xp} & $\Pi^P_2$-complete & \checkmark & Local & \checkmark & 
%         Focuses on coverage and irredundancy. Very high complexity for exact computation. \\
        
%         \bottomrule
%     \end{tabularx}
% \end{table}



